\section{Davis--Kahan Failure-Mode Sweep}
\label{sec:failure}

UET Theorem~4.1 \citep{uet2026} (subspace alignment) states that if the
embedding covariance $\Sigma$ has an eigen-gap of at least $2$ between
the top-$k$ and remaining eigenvalues, and $k < d$, then the leading
eigenspace aligns with the causal subspace up to a Davis--Kahan bound
\citep{davis1970rotation,yu2015useful}:
\[
\sin\theta(\hat U_k, U_k^\star) \;\leq\;
\frac{2\|\Sigma - \Sigma^\star\|_{\mathrm{op}}}{\lambda_k - \lambda_{k+1}}.
\]
The theorem is conditional on (i) a sufficient eigen-gap and (ii) $k < d$.
We sweep both conditions to quantify \emph{how much} each matters in
practice.

\subsection{Sweep design}
We generate $d \in \{16, 32, 64, 128, 256, 512\}$,
$k \in \{2, 4, 8, 16, 32, 64, 128\}$,
$\mathrm{gap\_ratio} \in \{0.5, 1, 2, 5, 10\}$. Each cell is averaged
over $10$ independent noise realisations
($2{,}100$ rows total). For each row we compute:
\begin{itemize}[leftmargin=1.5em,itemsep=1pt]
\item $\sin\theta$ between recovered top-$k$ and planted subspaces,
\item the Davis--Kahan theorem bound,
\item $\deff$ of the sample covariance,
\item a label \texttt{condition\_violated} $\in$
      \{\texttt{none}, \texttt{gap<2}, \texttt{k$\geq$d}\}.
\end{itemize}

\subsection{Headline numbers}

\begin{table}[H]
\centering
\small
\begin{tabular}{lSS}
\toprule
Regime & {$\overline{\sin\theta}$} & {fraction of cells} \\
\midrule
Condition satisfied (gap $\geq 2$, $k < d$) & 0.176 & 0.457 \\
Condition violated                          & 0.752 & 0.543 \\
Whole sweep                                 & 0.489 & 1.000 \\
\bottomrule
\end{tabular}
\caption{Mean subspace-alignment error $\sin\theta$ by regime. When the
theorem's sufficient condition holds, the recovered subspace aligns with
the planted subspace; when it fails (too small gap or $k \geq d$) the
alignment error jumps by a factor of $\mathbf{4.3\times}$.}
\label{tab:failure-headline}
\end{table}

\subsection{How $\sin\theta$ depends on the spectral gap}

\begin{figure}[H]
\centering
\begin{tikzpicture}
\begin{axis}[
  width=0.85\textwidth, height=6cm,
  xmode=log, xmin=0.4, xmax=12,
  xlabel={spectral gap ratio $\lambda_k / \lambda_{k+1}$},
  ylabel={$\sin\theta(\hat U_k, U_k^\star)$},
  grid=both, grid style={line width=0.1pt, draw=gray!30},
  legend style={font=\scriptsize, draw=none, at={(0.98,0.98)}, anchor=north east},
  tick label style={font=\small},
  label style={font=\small},
  every axis plot/.append style={thick},
]
  \addplot+[mark=*, color=blue!70!black]
    table[x=gap_ratio, y=sin_angle] {data/failure_k8_d16.dat};
  \addlegendentry{$d=16$}
  \addplot+[mark=square*, color=red!70!black]
    table[x=gap_ratio, y=sin_angle] {data/failure_k8_d64.dat};
  \addlegendentry{$d=64$}
  \addplot+[mark=triangle*, color=green!50!black]
    table[x=gap_ratio, y=sin_angle] {data/failure_k8_d256.dat};
  \addlegendentry{$d=256$}
  \addplot+[mark=diamond*, color=purple!70!black]
    table[x=gap_ratio, y=sin_angle] {data/failure_k8_d512.dat};
  \addlegendentry{$d=512$}

  \draw[dashed, thick, gray] (axis cs:2, 0) -- (axis cs:2, 0.35)
    node[pos=0.95, right, font=\scriptsize]{theorem boundary};
\end{axis}
\end{tikzpicture}
\caption{Subspace-alignment error vs. spectral-gap ratio, fixed $k=8$, for
four embedding widths. Left of the dashed line the theorem's
sufficient condition is violated and $\sin\theta$ is large and
$d$-independent; right of it $\sin\theta$ decays monotonically and the
curves almost overlap, confirming that the Davis--Kahan bound becomes
tight as the gap widens.}
\label{fig:failure-gap}
\end{figure}

\subsection{How violation frequency depends on $d$}

\begin{figure}[H]
\centering
\begin{tikzpicture}
\begin{axis}[
  width=0.70\textwidth, height=5cm,
  ybar, bar width=8pt,
  xlabel={embedding dim $d$}, xmode=log, log basis x=2,
  ylabel={fraction of cells with violation},
  ymin=0, ymax=1,
  grid=both, grid style={line width=0.1pt, draw=gray!25},
  xtick={16,32,64,128,256,512},
  xticklabels={16,32,64,128,256,512},
  tick label style={font=\small},
  label style={font=\small},
]
  \addplot+[fill=red!50!orange, draw=black]
    table[x=d, y=violated] {data/failure_by_d.dat};
\end{axis}
\end{tikzpicture}
\caption{Fraction of swept cells violating the theorem's sufficient
condition as a function of $d$. Because we keep $k$ fixed up to $128$,
the fraction drops sharply once $d$ exceeds the $k$-range, in line with
the $k < d$ clause of the theorem.}
\label{fig:failure-d}
\end{figure}

\begin{remark}
Figures~\ref{fig:failure-gap}--\ref{fig:failure-d} together form a
two-sided falsification test: \emph{both} the gap clause \emph{and} the
$k < d$ clause are necessary. Removing either produces large alignment
error. This is a stronger empirical statement than ``the theorem
holds'', because it shows the theorem is also \emph{tight} --- the
regime where it fails is empirically non-trivial.
\end{remark}
